\documentclass[a4paper,12pt]{article}
\usepackage[utf8]{inputenc}
\usepackage[polish]{babel}
\usepackage{geometry}
\geometry{a4paper, margin=1in}

\title{Raport z aktualnego stanu prac nad projektem}
\author{}
\date{\today}

\begin{document}

\maketitle

\section*{1. Aktualny stan prac}
\begin{itemize}
    \item \textbf{Symulacja demograficzna:} Została pomyślnie uruchomiona i przeprowadzona dla populacji 50,000 agentów na przestrzeni 50 lat. Wyniki symulacji obejmują:
    \begin{itemize}
        \item Piramidy wieku (statyczne i animowane),
        \item Trendy populacyjne,
        \item Dynamikę gospodarstw domowych,
        \item Rozkład płci w czasie.
    \end{itemize}
    \item \textbf{Optymalizacja parametrów:} Przeprowadzono proces GridSearch w celu znalezienia optymalnych parametrów modelu (np. wskaźników urodzeń, śmiertelności, migracji). Wyniki optymalizacji zostały uwzględnione w symulacji.
    \item \textbf{Dane wejściowe:} Dane demograficzne zostały przygotowane w oparciu o realistyczne struktury populacyjne inspirowane danymi GUS. Uwzględniono:
    \begin{itemize}
        \item Tabele śmiertelności i płodności,
        \item Rozkład wiekowy i płciowy,
        \item Czynniki ryzyka i choroby.
    \end{itemize}
\end{itemize}

\section*{2. Gotowość do etapu modelowania}
\begin{itemize}
    \item Dane są w pełni przygotowane do rozpoczęcia etapu modelowania. Symulacja dostarcza szczegółowych wyników, które mogą być wykorzystane do dalszej analizy i modelowania.
    \item Wyniki są dostępne w formie interaktywnych wizualizacji oraz szczegółowych statystyk, co ułatwia ich interpretację.
\end{itemize}

\section*{3. Ewentualne trudności i przeszkody}
\begin{itemize}
    \item \textbf{Brak problemów:} Wszystkie etapy prac przebiegły bez zakłóceń. Nie napotkano trudności technicznych ani merytorycznych, które mogłyby wpłynąć na realizację projektu.
    \item \textbf{Uwagi:} Wprowadzone zmiany w modelu (np. zwiększenie śmiertelności kobiet) zostały przetestowane i nie wpłynęły negatywnie na stabilność symulacji.
\end{itemize}

\section*{4. Podsumowanie}
Projekt jest na dobrej drodze do realizacji. Dane są gotowe do dalszego etapu modelowania, a dotychczasowe wyniki potwierdzają poprawność założeń i implementacji. W razie potrzeby możliwe jest dalsze dostosowanie parametrów modelu lub rozszerzenie analizy.

\vspace{1cm}
\noindent Jeśli potrzebne są dodatkowe informacje lub szczegółowe wyniki, proszę o kontakt.

\end{document}